\section{Micro-operações do Conjunto de Instruções Utilizado}

\label{sec:instrucoes}

\subsection{Convenções adotadas}

Para cada instrução, apresentamos as micro-operações dos três ciclos
pedidos (busca, indireto e execução) usando notação de transferência entre
registradores (RTL). Vale estabelecer as seguintes convenções usadas, algumas delas
são evidenciadas aqui pois não serão repetidas novamente ao decorrer do relatório:

\begin{description}
\item[\texttt{PC} (\textit{Program Counter}) -] Registrador que armazena o endereço da próxima instrução a ser executada.

\item[\texttt{IR} (\textit{Instruction Register}) -] Registrador que armazena a instrução atualmente em execução.

\item[\texttt{Reg[x]} -] Conteúdo do registrador de propósito geral \texttt{x}.

\item[\texttt{MEM[end]} -] Conteúdo da posição de memória localizada no endereço \texttt{end}.

\item[\texttt{HI}, \texttt{LO} -] Registradores especiais utilizados para armazenar, respectivamente, a parte alta e a parte baixa dos resultados das instruções de multiplicação e divisão.

\item[\textit{sign\_extend} -] Operação que estende um valor imediato de 16 para 32 bits preservando seu sinal.

\item[\textit{zero\_extend} -] Operação que estende um valor imediato de 16 para 32 bits preenchendo os bits mais significativos com zeros.
\end{description}

\textbf{Ciclo indireto.} A arquitetura MIPS não possui modo de
endereçamento indireto, isto é, nenhuma instrução do conjunto usa um endereço
armazenado em memória para localizar seu operando. Mesmo as instruções de
acesso à memória (\texttt{lw}, \texttt{sw}) calculam o endereço efetivo
diretamente durante a execução, somando um registrador base a um
deslocamento imediato. Isso é consistente com a caracterização do MIPS
como uma arquitetura \textit{load-store} feita no Apêndice A
\cite{patterson2017}, que descreve \texttt{c(rx)} (a soma de uma
constante ao conteúdo de um registrador base) como o único modo de
endereçamento de memória realmente oferecido pela máquina. Por esse
motivo, o ciclo indireto não se aplica a absolutamente nenhuma instrução
usada neste programa. \\

\textbf{Ciclo de execução.} A decodificação/leitura de operandos e a escrita do resultado
(\textit{write-back}) são apresentadas como parte do ciclo de execução de
cada instrução, em vez de fases isoladas. \\

\textbf{Repetição de ocorrências.} Quando a mesma instrução real aparece
mais de uma vez no código (inclusive quando gerada por diferentes
pseudo-instruções), a primeira ocorrência será detalhada por completo. No entanto, as
demais são apenas listadas por linha/rótulo. Isso se deve por executarem exatamente a mesma sequência de micro-operações.

\subsection{Pseudo-instruções e suas expansões}

A Tabela \ref{tab:pseudo} lista as pseudo-instruções adotadas e as instruções reais que o montador (MARS) gera a partir
delas. Essa distinção entre instruções reais e pseudo-instruções, assim
como a descrição de cada uma das pseudo-instruções usadas, é
documentada na Seção A.10 do Apêndice A \cite{patterson2017}, que trata
especificamente do assembly do MIPS R2000 simulado pelo SPIM. O
detalhamento de micro-operações de cada instrução real gerada está nas
subseções correspondentes, a seguir:

\begin{table}[h]
\centering
\small
\begin{tabularx}{\textwidth}{|l|l|X|}
\hline
\textbf{Pseudo-instrução} &
\textbf{Exemplo no código} &
\textbf{Instruções reais geradas} \\
\hline

\texttt{li \$reg, imed} &
\texttt{li \$v0, 4} (li\,14) &
\texttt{ori \$reg, \$zero, imed}$^{*}$ \\
\hline

\texttt{la \$reg, rótulo} &
\texttt{la \$a0, msg\_m} (li\,15) &
\texttt{lui \$reg, alta(rótulo)}; \texttt{ori \$reg, \$reg, baixa(rótulo)} \\
\hline

\texttt{lw \$reg, rótulo} &
\texttt{lw \$t0, m} (li\,68) &
\texttt{lui \$at, alta(rótulo)}; \texttt{lw \$reg, baixa(rótulo)(\$at)} \\
\hline

\texttt{sw \$reg, rótulo} &
\texttt{sw \$v0, m} (li\,21) &
\texttt{lui \$at, alta(rótulo)}; \texttt{sw \$reg, baixa(rótulo)(\$at)} \\
\hline

\texttt{move \$rd, \$rs} &
\texttt{move \$t0, \$v0} (li\,37) &
\texttt{addu \$rd, \$rs, \$zero} \\
\hline

\texttt{beqz \$rs, rótulo} &
\texttt{beqz \$v0, nao\_dividiu} (li\,104) &
\texttt{beq \$rs, \$zero, rótulo} \\
\hline

\texttt{bnez \$rs, rótulo} &
\texttt{bnez \$t1, testa\_n} (li\,156) &
\texttt{bne \$rs, \$zero, rótulo} \\
\hline

\texttt{mul \$rd, \$rs, \$rt} &
\texttt{mul \$s0, \$s0, \$s1} (li\,106) &
\texttt{mult \$rs, \$rt}; \texttt{mflo \$rd} \\
\hline

\end{tabularx}
\caption{Pseudo-instruções utilizadas e suas expansões}
\label{tab:pseudo}
\end{table}

Nas subseções seguintes, cada instrução real gerada é detalhada individualmente, na ordem de sua primeira aparição no código. Cada descrição segue uma estrutura comum, composta pelo formato da instrução, sua operação em nível de transferência entre registradores (RTL), sua utilização no contexto do programa e um exemplo correspondente à primeira ocorrência no código, seguido das demais ocorrências quando aplicável.

O campo \textbf{Formato (I, R ou J)} apresenta a organização dos 32 bits da instrução segundo a arquitetura MIPS, indicando a disposição de seus campos. Entre esses campos, o \texttt{opcode} identifica a operação a ser executada pelo processador e, nas instruções do formato R, o campo \texttt{funct} complementa essa identificação, distinguindo operações que compartilham o mesmo \texttt{opcode}. Os valores apresentados para esses campos seguem a especificação do manual de referência da arquitetura \cite{mips2016}.

Começamos pelas instruções geradas pela pseudo-instrução \texttt{li}, usada para carregar constantes pequenas em registradores.


% COMEÇO DA BAGAÇA TODA -------------------------------

\subsection{Instrução \texttt{ori}}

A instrução \texttt{ori} (\textit{OR Immediate}) realiza uma operação lógica
OR entre o conteúdo de um registrador e um valor imediato de 16 bits. Neste
programa, ela é utilizada principalmente na expansão da pseudo-instrução
\texttt{li}, permitindo carregar constantes pequenas diretamente em um
registrador.

\begin{instrucao}{ori}
\begin{description}
\item[Formato (I):] \texttt{opcode(6)=001101 | rs(5) | rt(5) | imediato(16)}
\item[Execução:] $Reg[rt] \leftarrow Reg[rs] \lor zero\_extend(imediato)$
\item[Uso:] Implementa a carga de constantes de \texttt{li} (com
$rs = \$zero$) e a metade menos significativa do endereço em \texttt{la}
e nas formas de \texttt{lw}/\texttt{sw} com rótulo direto.
\item[Primeira ocorrência:] \texttt{li \$v0, 4} (linha 14)
$\rightarrow$ \texttt{ori \$v0, \$zero, 4}:
$Reg[v0] \leftarrow Reg[zero] \lor zero\_extend(4) = 4$.
\item[Demais ocorrências:] linhas 15, 19, 24, 25, 29, 40, 41, 45, 50,
63, 64, 68, 69, 72, 95, 97, 149, 160, 171.
\end{description}
\end{instrucao}

\subsection{Instrução \texttt{lui}}

A instrução \texttt{lui} (\textit{Load Upper Immediate}) carrega um valor
imediato nos 16 bits mais significativos de um registrador, preenchendo os
16 bits inferiores com zeros. No programa, ela é empregada pelo montador na
expansão de pseudo-instruções que necessitam construir endereços completos
de memória.

\begin{instrucao}{lui}
\begin{description}
\item[Formato (I):] \texttt{opcode(6)=001111 | rs(5, não usado) | rt(5) | imediato(16)}
\item[Execução:] $Reg[rt] \leftarrow imediato \ll 16$
(bits menos significativos preenchidos com zero)
\item[Uso:] Usada como primeira metade de \texttt{la} e das formas de
\texttt{lw}/\texttt{sw} com rótulo direto, para carregar os 16 bits mais
significativos do endereço do rótulo.
\item[Primeira ocorrência:] primeira metade de
\texttt{la \$a0, msg\_m} (linha 15) $\rightarrow$
\texttt{lui \$a0, alta(msg\_m)}:
$Reg[a0] \leftarrow alta(msg\_m) \ll 16$.
\item[Demais ocorrências:] linhas 21, 25, 31, 41, 68, 69, 95, 97.
\end{description}
\end{instrucao}

\subsection{Instrução \texttt{syscall}}

Encerrada a carga de constantes e endereços, o programa realiza sua
primeira chamada de sistema, para exibir uma mensagem ao usuário. Os
códigos de serviço correspondem à
tabela de chamadas de sistema do SPIM (Seção A.9, Figura A.9.1,
\cite{patterson2017}), que também especifica a convenção de passagem de
parâmetros, no caso, o código do serviço em \texttt{\$v0} e os argumentos em
\texttt{\$a0}--\texttt{\$a3}.

\begin{instrucao}{syscall}
\begin{description}
\item[Formato (especial):] \texttt{opcode(6)=000000 | code(20) | funct(6)=001100}
\item[Execução:] dispara uma exceção (\textit{trap}) que transfere o
controle ao núcleo do simulador. O núcleo lê o código de serviço em
\texttt{\$v0} e executa a rotina correspondente (imprimir string usando o
endereço em \texttt{\$a0}, ler inteiro devolvendo-o em \texttt{\$v0},
imprimir inteiro usando \texttt{\$a0}, ou encerrar o programa), retornando
o controle à instrução seguinte.
\item[Primeira ocorrência:] linha 16, com $Reg[v0] = 4$: imprime a
string apontada por \texttt{\$a0} (mensagem "Digite m: ").
\item[Demais ocorrências:] linhas 20, 26, 30, 42, 47, 51.
\end{description}
\end{instrucao}

\subsection{Instrução \texttt{sw}}

A instrução \texttt{sw} (\textit{Store Word}) armazena uma palavra de 32 bits
de um registrador em uma posição de memória. No programa, ela é utilizada
para atualizar as variáveis globais após a execução das operações de cálculo.

\begin{instrucao}{sw}
\begin{description}
\item[Formato (I):] \texttt{opcode(6)=101011 | rs(5) | rt(5) | deslocamento(16)}
\item[Execução:] $EA \leftarrow Reg[rs] + sign\_extend(deslocamento)$;
$MEM[EA] \leftarrow Reg[rt]$
\item[Primeira ocorrência:] \texttt{sw \$v0, m} (linha 21), expandida
para \texttt{sw \$v0, baixa(m)(\$at)}:
$EA \leftarrow Reg[at] + sign\_extend(baixa(m))$;
$MEM[EA] \leftarrow Reg[v0]$
(guarda o valor de \texttt{m} lido do teclado).
\item[Demais ocorrências:] linhas 31, 59, 60, 61, 96, 98, 99, 140, 141,
142, 143, 159, 170.
\end{description}
\end{instrucao}

\subsection{Instrução \texttt{jal}}

A instrução \texttt{jal} (\textit{Jump And Link}) realiza a chamada de uma
sub-rotina, desviando a execução e armazenando o endereço de retorno no
registrador \texttt{\$ra}. No programa, ela é utilizada para invocar as
sub-rotinas responsáveis pelo cálculo do MMC e pelas divisões sucessivas.

\begin{instrucao}{jal}
\begin{description}
\item[Formato (J):] \texttt{opcode(6)=000011 | endereço-alvo(26)}
\item[Execução:] $Reg[ra] \leftarrow PC$
(endereço já incrementado na busca, ou seja, o da instrução seguinte);
$PC \leftarrow PC[31{:}28] \,\|\, (\text{endereço-alvo} \ll 2)$
\item[Primeira ocorrência:] \texttt{jal mmc} (linha 34): salva em
\texttt{\$ra} o endereço de retorno para \texttt{main} e desvia para o
rótulo \texttt{mmc}.
\item[Demais ocorrências:] linha 101 (\texttt{jal divide}).
\end{description}
\end{instrucao}

\subsection{Instrução \texttt{addu}}

A instrução \texttt{addu} (\textit{Add Unsigned}) realiza a soma entre dois
registradores sem considerar detecção de overflow. Neste programa, ela é
utilizada pelo montador para implementar a pseudo-instrução \texttt{move},
copiando o conteúdo de um registrador para outro.

\begin{instrucao}{addu}
\begin{description}
\item[Formato (R):] \texttt{opcode(6)=000000 | rs(5) | rt(5) | rd(5) | shamt(5) | funct(6)=100001}
\item[Execução:] $Reg[rd] \leftarrow Reg[rs] + Reg[rt]$
(soma sem verificação de \textit{overflow})
\item[Uso:] Implementa \texttt{move}, sempre com $rt = \$zero$, de modo
que $Reg[rd] \leftarrow Reg[rs]$.
\item[Primeira ocorrência:] \texttt{move \$t0, \$v0} (linha 37)
$\rightarrow$ \texttt{addu \$t0, \$v0, \$zero}:
$Reg[t0] \leftarrow Reg[v0] + 0$
(copia o MMC retornado por \texttt{mmc} para um registrador temporário).
\item[Demais ocorrências:] linhas 46, 115.
\end{description}
\end{instrucao}

\subsection{Instrução \texttt{addi}}

A instrução \texttt{addi} (\textit{Add Immediate}) soma um valor imediato ao
conteúdo de um registrador, armazenando o resultado em um registrador de
destino. No programa, ela é empregada principalmente para atualizar
contadores, ajustar ponteiros de pilha e incrementar variáveis de controle.

\begin{instrucao}{addi}
\begin{description}
\item[Formato (I):] \texttt{opcode(6)=001000 | rs(5) | rt(5) | imediato(16)}
\item[Execução:] $Reg[rt] \leftarrow Reg[rs] + sign\_extend(imediato)$
\item[Uso:] Diferente de \texttt{ori}, o imediato aqui é estendido com
sinal, e (pela especificação do MIPS) a instrução pode gerar exceção de
\textit{overflow} -- o que não ocorre neste programa, já que os valores
manipulados são pequenos.
\item[Primeira ocorrência:] \texttt{addi \$sp, \$sp, -12}
(linha 58): $Reg[sp] \leftarrow Reg[sp] + (-12)$,
abrindo espaço na pilha para o prólogo de \texttt{mmc}.
\item[Demais ocorrências:] linhas 94, 103, 110, 121, 139, 178.
\end{description}
\end{instrucao}

\subsection{Instrução \texttt{lw}}

A instrução \texttt{lw} (\textit{Load Word}) realiza a leitura de uma palavra
de 32 bits da memória para um registrador de propósito geral. Neste programa,
ela é utilizada para carregar os valores das variáveis armazenadas em memória
antes de sua manipulação pelas sub-rotinas.

\begin{instrucao}{lw}
\begin{description}
\item[Formato (I):] \texttt{opcode(6)=100011 | rs(5) | rt(5) | deslocamento(16)}
\item[Execução:] $EA \leftarrow Reg[rs] + sign\_extend(deslocamento)$;
$Reg[rt] \leftarrow MEM[EA]$
\item[Primeira ocorrência:] \texttt{lw \$t0, m} (linha 68), expandida
para \texttt{lw \$t0, baixa(m)(\$at)}:
$EA \leftarrow Reg[at] + sign\_extend(baixa(m))$;
$Reg[t0] \leftarrow MEM[EA]$
(carrega o valor atual de \texttt{m}).
\item[Demais ocorrências:] linhas 69, 118, 119, 120, 145, 146, 147, 152,
163, 174, 175, 176, 177.
\end{description}
\end{instrucao}

\subsection{Instrução \texttt{bne}}

A instrução \texttt{bne} (\textit{Branch if Not Equal}) realiza um desvio
condicional quando os conteúdos de dois registradores são diferentes. No
programa, ela é empregada para controlar laços de repetição e implementar a
pseudo-instrução \texttt{bnez}.

\begin{instrucao}{bne}
\begin{description}
\item[Formato (I):] \texttt{opcode(6)=000101 | rs(5) | rt(5) | deslocamento(16)}
\item[Execução:] se $Reg[rs] \neq Reg[rt]$, então
$PC \leftarrow PC + (sign\_extend(deslocamento) \ll 2)$;
caso contrário $PC$ permanece o valor já incrementado na busca.
\item[Primeira ocorrência:]
\texttt{bne \$t0, \$t2, tentaDividir} (linha 75):
desvia para \texttt{tentaDividir} se \texttt{m} (em \texttt{\$t0}) ainda
for diferente de 1.
\item[Demais ocorrências:] linha 81 (\texttt{bne} direto); linhas 156 e
167 (\texttt{bnez} expandido, com $rt = \$zero$).
\end{description}
\end{instrucao}

\subsection{Instrução \texttt{j}}

A instrução \texttt{j} (\textit{Jump}) realiza um salto incondicional para um
novo endereço de execução. Neste programa, ela é utilizada para alterar o
fluxo de controle entre diferentes trechos do algoritmo, sem necessidade de
armazenar um endereço de retorno.

\begin{instrucao}{j}
\begin{description}
\item[Formato (J):] \texttt{opcode(6)=000010 | endereço-alvo(26)}
\item[Execução:] $PC \leftarrow PC[31{:}28] \,\|\, (\text{endereço-alvo} \ll 2)$
\item[Primeira ocorrência:] \texttt{j fim\_mmc} (linha 83): desvia
incondicionalmente para o epílogo de \texttt{mmc}.
\item[Demais ocorrências:] linhas 107, 111.
\end{description}
\end{instrucao}

\subsection{Instrução \texttt{beq}}

A instrução \texttt{beq} (\textit{Branch if Equal}) realiza um desvio
condicional quando os conteúdos de dois registradores são iguais. Neste
programa, ela é utilizada na implementação de decisões lógicas e na expansão
da pseudo-instrução \texttt{beqz}.

\begin{instrucao}{beq}
\begin{description}
\item[Formato (I):] \texttt{opcode(6)=000100 | rs(5) | rt(5) | deslocamento(16)}
\item[Execução:] se $Reg[rs] = Reg[rt]$, então
$PC \leftarrow PC + (sign\_extend(deslocamento) \ll 2)$.
\item[Única ocorrência:] \texttt{beqz \$v0, nao\_dividiu}
(linha 104), expandida para \texttt{beq \$v0, \$zero, nao\_dividiu}:
desvia se \texttt{divide} não dividiu nenhum dos dois valores.
\end{description}
\end{instrucao}

\subsection{Instrução \texttt{mult}}

A instrução \texttt{mult} (\textit{Multiply}) realiza a multiplicação de dois
operandos armazenados em registradores, gravando o resultado de 64 bits nos
registradores especiais \texttt{HI} e \texttt{LO}. No programa, ela faz parte
da expansão da pseudo-instrução \texttt{mul}.

\begin{instrucao}{mult}
\begin{description}
\item[Formato (R especial):]
\texttt{opcode(6)=000000 | rs(5) | rt(5) | (não usado, 10) | funct(6)=011000}
-- não grava em registrador de propósito geral.
\item[Execução:] $\{HI, LO\} \leftarrow Reg[rs] \times Reg[rt]$
(produto de 64 bits com sinal; \texttt{HI} recebe a parte alta,
\texttt{LO} a parte baixa)
\item[Única ocorrência:] primeira metade de
\texttt{mul \$s0, \$s0, \$s1} (linha 106) $\rightarrow$
\texttt{mult \$s0, \$s1}:
$\{HI, LO\} \leftarrow Reg[s0] \times Reg[s1]$.
\end{description}
\end{instrucao}

\subsection{Instrução \texttt{mflo}}

A instrução \texttt{mflo} (\textit{Move From LO}) copia para um registrador de
propósito geral o conteúdo do registrador especial \texttt{LO}. Neste
programa, ela é utilizada para recuperar o resultado das multiplicações
realizadas pela instrução \texttt{mult}.

\begin{instrucao}{mflo}
\begin{description}
\item[Formato (R especial):]
\texttt{opcode(6)=000000 | (não usado, 10) | rd(5) | (não usado, 5) |
funct(6)=010010}
\item[Execução:] $Reg[rd] \leftarrow LO$
\item[Primeira ocorrência:] segunda metade de
\texttt{mul \$s0, \$s0, \$s1} (linha 106) $\rightarrow$
\texttt{mflo \$s0}:
$Reg[s0] \leftarrow LO$
(acumula o divisor no MMC).
\item[Demais ocorrências:] linhas 158, 169
(quociente da divisão em \texttt{divide}).
\end{description}
\end{instrucao}

\subsection{Instrução \texttt{jr}}

A instrução \texttt{jr} (\textit{Jump Register}) transfere o controle para o
endereço contido em um registrador. Neste programa, ela é utilizada para
retornar de sub-rotinas por meio do endereço previamente armazenado em
\texttt{\$ra}.

\begin{instrucao}{jr}
\begin{description}
\item[Formato (R):] \texttt{opcode(6)=000000 | rs(5) | (não usado, 15) | funct(6)=001000}
\item[Execução:] $PC \leftarrow Reg[rs]$
\item[Primeira ocorrência:] \texttt{jr \$ra} (linha 122): retorna de
\texttt{mmc} para \texttt{main}, usando o endereço restaurado da pilha.
\item[Demais ocorrências:] linha 179 (retorno de \texttt{divide}).
\end{description}
\end{instrucao}

\subsection{Instrução \texttt{div}}

A instrução \texttt{div} (\textit{Divide}) realiza a divisão inteira entre
dois registradores, armazenando o quociente em \texttt{LO} e o resto em
\texttt{HI}. No programa, ela é utilizada para verificar se um divisor
candidato divide exatamente os valores utilizados no cálculo do MMC.

\begin{instrucao}{div}
\begin{description}
\item[Formato (R especial):] \texttt{opcode(6)=000000 | rs(5) | rt(5) | (não usado, 10) | funct(6)=011010}
\item[Execução:] $LO \leftarrow Reg[rs] \div Reg[rt]$
(quociente, com sinal);
$HI \leftarrow Reg[rs] \bmod Reg[rt]$
(resto)
\item[Primeira ocorrência:] \texttt{div \$t0, \$a2} (linha 153):
$LO \leftarrow Reg[t0] \div Reg[a2]$;
$HI \leftarrow Reg[t0] \bmod Reg[a2]$
(testa se \texttt{*m} é divisível por \texttt{d}).
\item[Demais ocorrências:] linha 164 (mesmo teste para \texttt{*n}).
\end{description}
\end{instrucao}

\subsection{Instrução \texttt{mfhi}}

A instrução \texttt{mfhi} (\textit{Move From HI}) copia para um registrador de
propósito geral o conteúdo do registrador especial \texttt{HI}. Neste
programa, ela é utilizada para verificar o resto produzido pela instrução
\texttt{div}, permitindo determinar se uma divisão foi exata.

\begin{instrucao}{mfhi}
\begin{description}
\item[Formato (R especial):] \texttt{opcode(6)=000000 | (não usado, 10) | rd(5) | (não usado, 5) | funct(6)=010000}
\item[Execução:] $Reg[rd] \leftarrow HI$
\item[Primeira ocorrência:] \texttt{mfhi \$t1} (linha 154):
$Reg[t1] \leftarrow HI$
(recupera o resto da divisão de \texttt{*m} por \texttt{d}).
\item[Demais ocorrências:] linha 165 (resto da divisão de \texttt{*n}).
\end{description}
\end{instrucao}